% Part: many-valued-logic
% Chapter: three-valued-logics
% Section: goedel

\documentclass[../../../include/open-logic-section]{subfiles}

\begin{document}

\olfileid[fr]{mvl}{thr}{god}

\olsection{Les logiques de G\"odel}

Kurt G\"odel a introduit une suite
de logiques à $n$ valeurs. Toutes les
formules valides en logique
intuitionniste y sont valides, et
toutes les formules qui y sont valides
le sont en logique classique. Voici
le premier cas intéressant.

\begin{defn}\ollabel{defn:goedel}
La \emph{logique de G\"odel à trois
valeurs}~$\LogGod$ est définie par
la matrice suivante :
\begin{enumerate}
  \item Le langage propositionnel usuel
  $\Lang L_0$ avec $\lfalse$,
  $\lnot$, $\land$, $\lor$ et $\lif$.
  \item L'ensemble des valeurs de vérité
  $V=\{\True,\Undef,\False\}$.
  \item $\True$ est la seule valeur
  désignée, c'est-à-dire
  $V^+=\{\True\}$.
  \item Pour $\lfalse$, on a
  $\tf{\lfalse}=\False$. Les fonctions
  de vérité des autres connecteurs
  sont données par les tables suivantes :
  \begin{center}
    \begin{tabular}{c|c}
      $\tf{\lnot}[\LogGod]$ & \\
      \hline
      $\True$ & $\False$ \\
      $\Undef$ & $\False$ \\
      $\False$ & $\True$
    \end{tabular}
    \quad
    \begin{tabular}{c|ccc}
      $\tf{\land}[\LogGod]$ & $\True$ & $\Undef$ & $\False$ \\
      \hline
      $\True$ & $\True$ & $\Undef$ & $\False$ \\
      $\Undef$ & $\Undef$ & $\Undef$ & $\False$ \\
      $\False$ & $\False$ & $\False$ & $\False$
    \end{tabular}
    \\[2ex]
    \begin{tabular}{c|ccc}
      $\tf{\lor}[\LogGod]$ & $\True$ & $\Undef$ & $\False$ \\
      \hline
      $\True$ & $\True$ & $\True$ & $\True$ \\
      $\Undef$ & $\True$ & $\Undef$ & $\Undef$ \\
      $\False$ & $\True$ & $\Undef$ & $\False$
    \end{tabular}
    \quad
    \begin{tabular}{c|ccc}
      $\tf{\lif}[\LogGod]$ & $\True$ & $\Undef$ & $\False$ \\
      \hline
      $\True$ & $\True$ & $\Undef$ & $\False$ \\
      $\Undef$ & $\True$ & $\True$ & $\False$ \\
      $\False$ & $\True$ & $\True$ & $\True$
    \end{tabular}
  \end{center}
\end{enumerate}
\end{defn}

Les tables de $\land$ et de~$\lor$ sont communes à la logique
de G\"odel, à la logique de \L ukasiewicz et à la logique forte de
Kleene. Les tables de $\lnot$ et de~$\lif$ de la logique de
G\"odel diffèrent de celles des deux autres logiques.
Dans la logique de G\"odel,
$\tf{\lnot}(\Undef)=\False$.
Contrairement aux logiques de
\L ukasiewicz et de Kleene,
$\tf{\lif}(\Undef,\False)=\False$;
contrairement à la logique de Kleene,
mais comme dans celle de \L ukasiewicz,
$\tf{\lif}(\Undef,\Undef)=\True$.

Comme le suggère le lien évoqué plus
haut, $\LogGod[3]$ est proche de la
logique intuitionniste. Toutes les
formules valides intuitionnistement
sont des tautologies de~$\LogGod[3]$;
en revanche, nombre de tautologies
classiques qui ne sont pas valides
intuitionnistement ne sont pas non
plus des tautologies de~$\LogGod[3]$.
Par exemple, les formules suivantes
ne sont pas des tautologies :
\begin{align*}
  & p\lor\lnot p && (p\lif q)\lif(\lnot p\lor q) \\
  & \lnot\lnot p\lif p && \lnot(\lnot p\land\lnot q)\lif(p\lor q) \\
  & ((p\lif q)\lif p)\lif p && \lnot(p\lif q)\lif(p\land\lnot q)
\end{align*}
Inversement, certaines tautologies
de~$\LogGod[3]$ ne sont pas valides
intuitionnistement, par exemple
$\lnot\lnot p\lor\lnot p$ ou
$(p\lif q)\lor(q\lif p)$.

\begin{prob}
  Donner des tables de vérité montrant
  que les formules suivantes sont des
  tautologies de~$\LogGod[3]$ :
  \begin{align*}
    & \lnot\lnot p\lor\lnot p \\
    & (p\lif q)\lor(q\lif p) \\
    & \lnot(p\land q)\lif(\lnot p\lor\lnot q) \\
    & (p\lif q)\lor(q\lif r)\lor(r\lif s)
  \end{align*}
\end{prob}

\begin{prob}
  Donner des tables de vérité montrant
  que les formules suivantes ne sont
  pas des tautologies de~$\LogGod[3]$ :
  \begin{align*}
    & (p\lif q)\lif(\lnot p\lor q) \\
    & \lnot(\lnot p\land\lnot q)\lif(p\lor q) \\
    & ((p\lif q)\lif p)\lif p \\
    & \lnot(p\lif q)\lif(p\land\lnot q)
  \end{align*}
\end{prob}

\begin{prob}
  Parmi les relations suivantes,
  lesquelles valent dans la logique
  de G\"odel? Donner une table de
  vérité pour chacune.
  \begin{enumerate}
    \item $p,p\lif q\Entails q$;
    \item $p\lor q,\lnot p\Entails q$;
    \item $p\land q\Entails p$;
    \item $p\Entails p\land p$;
    \item $p\Entails p\lor q$.
  \end{enumerate}
\end{prob}

\end{document}
